\section{Certificate tails and numerical dense bounds}\label{sec:certificate}
The factorial proof is complete. This independent route keeps information that its central-block comparison does not: explicit bounds on the exponential density cost at fixed tracked-card densities. We first count fixed certificates for many short cycles, then estimate reciprocal cycle lengths and pass to all representation shapes. Here $n=2^d$ is the full-deck size and $r$ is the number of tracked cards. Thus the exponent in \eqref{net:exponent} is used with $N=n$ and $k=r$. Logarithms in this section are to base two unless $\ln$ is written.

\subsection{A fixed-certificate cycle tail}\label{net:certificate}
\begin{proposition}[High-height cycle tail]\label{net:hightail}
There are absolute $b,\delta>0$ and an integer $J_0\ge1$ such that, uniformly in the tracked input injection of size $r$ and for every integer $J\ge J_0$,
\begin{align}
 \log_2\E 2^{b\sum_{j\ge J}C_j}&\le O(r2^{-bJ}),\label{net:highmgf}\\
 \log_2\E 2^{b\sum_jC_j}&\le O(r(r/n)^\delta).\label{net:sparsemgf}
\end{align}
The first bound also holds with expectation conditional on all input-side
switches and all output switches of heights at most $s$, for every integer
$0\le s\le d$ with $2s<J$, uniformly in the fixed switch values and with
the same absolute implied constant.
\end{proposition}
\begin{proof}
The estimates are immediate for $r=0$, so assume $r\ge1$.
Let $\mathcal F_t$ be the sigma-field generated by all input-side
switches and the output switches of heights at most $t$.
The arrivals at every node are determined by the input switches.
At a height-$j$ node the two child maps determine its cycles; the
outer output switch only exchanges the two positions in an output
pair and does not change their pairing.  Thus $C_j$ is
$\mathcal F_{j-1}$-measurable.

Fix a nonnegative integer $t$ with $2t<j$ and condition on $\mathcal F_t$.
In each height-$j$ node expose one child's remaining switches.
In the other child the fixed low layers are followed by independent
parallel output butterflies of height $j-1-t$.  Since
$j-1-t\ge\lfloor j/2\rfloor$, at least
$m=\lfloor j/2\rfloor-1$ final output layers remain fresh.
This includes $t=(j-1)/2$ when $j$ is odd.
The two child bijections identify alternating cycles with cycles of
their relative permutation. Explore such a cycle through the fresh
child and back through the known inverse of the other. Its closing
target is specified, and the closing route through its remaining
butterfly is unique.

Suppose $C_j\ge s\ge r2^{-\eta j}$, where $\eta>0$ is a sufficiently small absolute constant. At least $s/2$ of some $s$ counted cycles have length at most $2r/s$. Among these choose $u=\max(1,\lfloor s/100\rfloor)$ cycles uniformly without replacement and choose a uniform cyclic start on each. A fresh-layer switch of a closing path meets at most one path of another cycle. The conditional chance of selecting that other short cycle is at most $(u-1)/(s/2-1)$, less than $1/49$ when $u>1$, and zero when $u=1$. Here we bound interference against all other selected cycles, including those later in the canonical exploration order; restricting to already explored cycles only decreases it. Thus mean cross-cycle interference on a closing path is less than $j/30$. Within its own cycle, Lemma~\ref{net:encounters} bounds the mean number of earlier encountered switches by $\log_2(2r/s)\le1+\eta j$. For large $j$, the expected total previously exposed switches on all closing paths is less than $uj/4$.

Consequently each realized network in this event admits a certificate of starts and lengths with at most $uj/4$ previously exposed closing switches. This is an existence assertion before taking a union bound; we do not condition the coin law on the chosen favorable cycles. There are at most
\[
 \binom ru(2r/s)^u
\]
fixed certificates. Order their starts canonically and explore each fixed certificate in that order. Define its success event to require that all prescribed cycles close and that at least $K=\lceil u(j/4-2)\rceil$ closing queries are fresh. The preceding low-interference existence argument supplies a certificate in this event for every network being counted. At every fresh closing query its required bit is already determined, so the query succeeds with conditional probability $1/2$, irrespective of intervening nonclosing queries. Abort if a closing bit fails or the certificate becomes invalid, and stop on the $K$th successful fresh closing query. The probability of reaching this stopping threshold without failure is at most $2^{-K}\le2^{-u(j/4-2)}$. This bounds the stated success event even when unsuccessful explorations make fewer than $K$ fresh queries. The base-two logarithm of the certificate count is at most
\[
 u\{\log_2(200er/s)+\log_2(2r/s)\}\le u(O(1)+2\eta j).
\]
Since $u\ge s/200$, reducing $\eta$ gives
$\PP(C_j\ge s\mid\mathcal F_t)\le2^{-cjs}$, uniformly in the
conditioning whenever $2t<j$.

Split the moment tail at $r2^{-\eta j}$. The deterministic part costs
$O(jr2^{-\eta j})$ and the remaining geometric tail is summable at
any sufficiently small fixed exponent. Decreasing that exponent
gives absolute $b',A>0$ and $J_0$ such that, for every $j\ge J_0$
and $2t<j$,
\begin{equation}\label{net:single-height-conditional}
 \log_2\E\bigl[2^{b'jC_j}\mid\mathcal F_t\bigr]
 \le Ar2^{-b'j}.
\end{equation}

We retain the conditioning while summing over heights. Fix integers
$J\ge J_0$ and $0\le s\le d$ with $2s<J$. If $J>d$, the
sum is empty. Otherwise put $L_q=2^qJ$ and, for each $q\ge0$
with $L_q\le d$, define
\begin{gather*}
 I_q=\{L_q,\ldots,\min(2L_q-1,d)\},\qquad
 B_q=\sum_{j\in I_q}C_j,\\
 \ell_q=|I_q|,\qquad
 t_q=\left\lfloor\frac{L_q-1}{2}\right\rfloor.
\end{gather*}
These bands partition the integer heights from $J$ to $d$.
We have $s\le t_q$, $2t_q<L_q$, and $\ell_q\le L_q$.
Choose $0<b\le b'/2$. Conditional H\"older within a band and
$2b\ell_q\le b'j$ for every $j\in I_q$ give
\begin{align*}
 \log_2\E\bigl[2^{2bB_q}\mid\mathcal F_{t_q}\bigr]
 &\le\frac1{\ell_q}\sum_{j\in I_q}
       \log_2\E\bigl[2^{2b\ell_q C_j}\mid\mathcal F_{t_q}\bigr]\\
 &\le\frac{Ar}{\ell_q}\sum_{j\in I_q}2^{-b'j}
 \le Ar2^{-b'L_q}.
\end{align*}

Split the band indices into their two parity classes, and let
$V_q=\sum_{p<q,\ p\equiv q\ (2)}B_p$.
For $q\ge2$, every earlier band in the same class ends at or
below $L_q/2-1\le t_q$. Since $C_j$ is
$\mathcal F_{j-1}$-measurable, $V_q$ is
$\mathcal F_{t_q}$-measurable; for $q=0,1$ it is zero.
Also $\mathcal F_s\subseteq\mathcal F_{t_q}$ for every band.
The conditional tower step is therefore
\begin{align*}
 \E\bigl[2^{2b(V_q+B_q)}\mid\mathcal F_s\bigr]
 &=\E\left[2^{2bV_q}
       \E\bigl[2^{2bB_q}\mid\mathcal F_{t_q}\bigr]
       \,\middle|\,\mathcal F_s\right]\\
 &\le 2^{Ar2^{-b'L_q}}
       \E\bigl[2^{2bV_q}\mid\mathcal F_s\bigr].
\end{align*}
Integrating the highest band first and then descending gives,
for $\varepsilon\in\{0,1\}$,
\[
 \E\left[2^{2b\sum_{q\equiv\varepsilon\ (2)}B_q}
       \,\middle|\,\mathcal F_s\right]
 \le 2^{Ar\sum_{q\equiv\varepsilon\ (2)}2^{-b'L_q}}.
\]
All band sums here are finite. Conditional Cauchy--Schwarz between
the two classes now yields
\[
 \log_2\E\left[2^{b\sum_{j\ge J}C_j}
       \,\middle|\,\mathcal F_s\right]
 \le\frac{Ar}{2}\sum_{q\ge0}2^{-b'2^qJ}
 \le A'r2^{-bJ}
\]
with an absolute $A'$. This proves the asserted conditional bound;
averaging it proves~\eqref{net:highmgf}. No independence between
different heights or bands has been assumed.

For \eqref{net:sparsemgf}, put $\rho=r/n$ and take $J=\lfloor .1\log_2(1/\rho)\rfloor$, or a fixed sufficiently large minimum. Below this height dominate cycles by shared input pairs. A law on occupancy bits is called strongly Rayleigh if its multiaffine generating polynomial has no zero when all variables have positive imaginary part. The deterministic occupancy law has a monomial generating polynomial. Fair transpositions preserve this property, and strongly Rayleigh laws are negatively associated: products of increasing functions on disjoint supports have expectation at most the product of their expectations~\cite[Theorems~4.9 and~4.20]{BorceaBrandenLiggett2009}. This is a fixed input-layer process, with no additional path conditioning. Indicators that both sites in a disjoint pair are marked are increasing functions on disjoint supports and inherit negative association. Their exponential moment is bounded by the product of the Bernoulli moments. A site's marginal before height $j$ is at most $2^j\rho$, and the two pre-switch input groups are independent because that coordinate has not yet been used. The expected number of shared pairs is therefore at most $r2^j\rho$. H\"older over the low heights introduces an exponent $O(J)$; decreasing the fixed $b$ if necessary, its cost is a positive power of $\rho$ times $r$.

The low-height dominating factor depends only on the input switches.
Conditioning on $\mathcal F_0$, pull that factor outside the
expectation and apply the uniform conditional high-height bound just
proved with $s=0$. Averaging over the input switches combines the
two costs and gives~\eqref{net:sparsemgf} for some absolute
$\delta>0$. This final estimate uses the input randomness.
If $\rho$ is bounded below, the fixed low heights cost $O(r)$,
which is consistent with~\eqref{net:sparsemgf} after changing the constant.
\end{proof}

\subsection{Dense tuple densities}\label{net:dense}
The sparse moment estimate alone does not provide a useful exponent at fixed density. Here a bound on reciprocal cycle lengths gives a smaller explicit density exponent. Define
\[
 h_j=\max\left\{\E(1/L):L\in\mathbb Z_{\ge1},\ \PP(L=l)\le l/2^{j-1}\right\},
\]
and
\[
 H(\rho)=\max\left\{\frac12\sum_{j\ge1}h_jx_j:
 0\le x_j\le1,\ \sum_{j\ge1}2^{-j}x_j\le\rho\right\}.
\]
\begin{proposition}[Dense truncation]\label{net:densetruncation}
For every $\rho_0\in(0,1]$ and $\epsilon>0$, there are
$c>0$ and $n_0$ such that, for every dyadic $n\ge n_0$ and
every integer $r$ with $\rho=r/n\in[\rho_0,1]$, the density
exponent \eqref{net:exponent} exceeds
\[
 r\left[H(\rho)+\frac{-\rho-(1-\rho)\ln(1-\rho)}{\rho\ln2}+\epsilon\right]
\]
with probability at most $\exp(-cr)$, uniformly in the input injection.
The constants depend only on $\rho_0,\epsilon$, and the expression
is interpreted continuously at $\rho=1$.
\end{proposition}
\begin{proof}
We first estimate the reciprocal length of a cycle in the relative
permutation of the two children at a height-$j$ node.  Put
$M=2^{j-1}$.  Fix one child's map and the other child's input-side
map, leaving its final output butterfly fresh.  Start at a uniform
boundary column.  For a proposed cycle length $l$, expose the first
$l-1$ edges.  The closing route is unique if compatible, and if $v$
of its switches have already been exposed, its conditional closing
probability is $2^v/M$.  Hence
\[
 \E[\ind_{\{L=l\}}2^{-v}]\le M^{-1}.
\]
For a fixed realized cycle, rotating its start makes the closing
path uniform among its $l$ paths.  Lemma~\ref{net:encounters}
bounds their shared switches by $l\log_2l/2$, so the mean of $v$
over starts is at most $\log_2l$.  Jensen's inequality gives a
mean weight at least $1/l$.  Averaging over configurations proves
$\PP(L=l)\le l/M$ for the uniform start.

Now average all the child-interior randomness.  Each child law is
invariant under conjugation by an XOR translation of its boundary
columns.  The relative permutation has the same invariance, and
these translations act transitively.  Therefore, for every specified
column $i$, its cycle length $L_i$ satisfies
$\E(1/L_i)\le h_j$.  This assertion uses the unconditioned child
interiors; it is not asserted after the maps fixed in the preceding
exploration have been revealed.

Let $D_v$ be the set of shared input columns at a height-$j$ node $v$:
both positions of such a column carry tracked cards.  Write $L_{v,i}$
for the full cycle length at that node and set
\[
 S_j=\sum_{v:\,\operatorname{height}(v)=j}|D_v|,
 \qquad Z_v=\sum_{i\in D_v}\frac1{L_{v,i}},
 \qquad Y_j=\sum_{v:\,\operatorname{height}(v)=j}Z_v.
\]
A tracked cycle uses every column of one full relative-permutation
cycle.  It contributes one to $Y_j$; cycles only partly contained
in $D_v$ add nonnegative terms.  Thus $C_j\le Y_j$.
Let $\mathcal A_j$ record the arrivals at this height and the
boundary input switches of its nodes.  These data depend only on
coins outside the child interiors.  The reciprocal-length estimate
therefore gives
\begin{equation}\label{net:dense-conditional-mean}
 C_j\le Y_j,\qquad
 \E(Y_j\mid\mathcal A_j)\le h_j S_j.
\end{equation}

The means of the shared counts satisfy the budget defining $H$.
Use the coordinate convention of Section~\ref{sec:factorial}, so
the input sweep visits $d,\ldots,1$.  Let $A_j$ count unordered
pairs of tracked original inputs that agree in coordinates
$1,\ldots,j-1$ and differ in coordinate $j$.  The least differing
coordinate assigns every pair to exactly one $A_j$.  Before
direction $j$, the two paths remain in different subcubes, so their
already updated coordinates are independent uniforms.  They meet
at direction $j$ with probability $2^{-(d-j)}$.  Consequently, with
$x_j=2\E S_j/r$,
\[
 \E S_j=A_j2^{-(d-j)},\qquad 0\le x_j\le1,\qquad
 \sum_j2^{-j}x_j
 =\frac{2\sum_jA_j}{rn}=\frac{r-1}{n}\le\rho.
\]
Here $S_j\le r/2$ gives $x_j\le1$.  Taking expectations
in~\eqref{net:dense-conditional-mean} now yields
$\sum_j h_j\E S_j\le rH(\rho)$.

We next turn this mean bound into a tail bound, without conditioning
several heights simultaneously.  Fix a large height cutoff $J$.
At each $j\le J$, the disjoint-pair occupancy indicators are
negatively associated by the fixed-input-layer argument above.
Their sum $S_j$, whose mean is at most $r/2$, has the Bernoulli-product
upper exponential moment.  For every fixed $t>0$, Chernoff's bound
therefore gives
\[
 \PP\{S_j>\E S_j+tr\}\le e^{-c_t r}.
\]
Conditional on $\mathcal A_j$, the variables $Z_v$ in different
nodes are independent, and $0\le Z_v\le2^{j-1}$.  There are
$n/2^j$ such nodes.  Hoeffding's inequality and
\eqref{net:dense-conditional-mean} give
\[
 \PP\{Y_j>h_jS_j+tr\mid\mathcal A_j\}
 \le \exp\{-c t^2r^2/(n2^j)\}
 \le \exp\{-c t^2\rho_0 r/2^j\}.
\]
The second inequality uses $r/n\ge\rho_0$.  A union bound over
the fixed set $j\le J$, taking $t$ small in terms of $J,\epsilon$,
controls their total excess over $\sum_{j\le J}h_j\E S_j$.
By~\eqref{net:highmgf}, choosing $J$ large enough in terms of
$\epsilon$ controls $\sum_{j>J}C_j$ at the remaining additive
slack with probability $e^{-c_\epsilon r}$.  Neither this choice
nor the low-height constants depend on the individual density
$\rho\in[\rho_0,1]$.  Enlarging $n_0$ absorbs the finite union
factor and shows that $\sum_jC_j$ exceeds
$rH(\rho)+\epsilon r/2$ with probability at most $e^{-cr}$.

Finally Stirling's formula, with $0!=1$ at the endpoint, gives
\[
 \log_2((n)_r/n^r)
 =r\frac{-\rho-(1-\rho)\ln(1-\rho)}{\rho\ln2}+O(\log n).
\]
The error is uniform for $1\le r\le n$.  Since $r\ge\rho_0n$,
it is at most $\epsilon r/2$ for $n\ge n_0(\rho_0,\epsilon)$.
Adding this deterministic normalization proves the proposition.
\end{proof}

The defining linear optimization fills the shortest allowed cycle lengths first. With $M=2^{j-1}$ and $q(q+1)\le2M<(q+1)(q+2)$,
\begin{equation}\label{net:hj}
 h_j=\frac qM+\frac{1-q(q+1)/(2M)}{q+1},
 \qquad h_j\le\sqrt{2/M}.
\end{equation}
The allocation ratios $2^{j-1}h_j$ increase with $j$. Indeed $Mh(M)$ is the maximum of $\sum_lz_l/l$ subject to $0\le z_l\le l$ and $\sum_lz_l=M$; increasing $M$ permits adding positive-valued mass. Thus the budget for $H$ is assigned at the highest indices first, with at most one partially filled index.

Here are explicit arithmetic bounds, with enough detail to reproduce the finite check. Formula~\eqref{net:hj} simplifies to $h_j=q/(2M)+1/(q+1)$. The inequality $h_j\le\sqrt{2/M}$ follows by putting $x=2M$: the convex quadratic $(q+x/(q+1))^2-4x$ is nonpositive at the two endpoints $q(q+1)$ and $(q+1)(q+2)$, hence throughout the interval. Therefore
\[
 \frac12\sum_{j>20}h_j\le\frac{1+\sqrt2}{1024}
 <\frac{169}{70\cdot1024},
 \qquad \sum_{j>20}2^{-j}=2^{-20}.
\]
Subtract this last budget, allocate the first twenty levels greedily, and add the rational value bound. Exact rational arithmetic gives
\[
 H(1/8)<.791402,\qquad H(13/50)<1.098069,
 \qquad H(1)<1.958069.
\]
These imply the stated useful bounds $.80$, $1.11$ and $1.97$.
For the logarithms in the entropy correction, use
\[
 -\ln(1-x)\le\sum_{i=1}^{80}\frac{x^i}{i}
       +\frac{x^{81}}{81(1-x)},\qquad
 \ln2=2\sum_{j=0}^{39}\frac{(1/3)^{2j+1}}{2j+1}+R,
 \quad 0<R\le\frac{2(1/3)^{81}}{81(1-1/9)}.
\]
The entropy numerator is negative, so increasing its positive denominator gives an upper bound. The resulting rational comparisons give a full bracket below $.904054$ at $\rho=7/25$ and below $.515374$ at $\rho=1$. In particular they imply the margins $.96$ and $.54$ used below.

These endpoint checks cover every $0<\rho\le7/25$. For $\rho\le1/8$, the entropy correction is nonpositive and $H(\rho)\le H(1/8)<.80$. On $[1/8,1/4]$ and $[1/4,7/25]$, the slopes of $H$ are respectively $7/3$ and $3/2$. The absolute derivative of the entropy correction is
\[
 \frac{-\rho-\ln(1-\rho)}{\rho^2\ln2}
 \le\frac1{2(1-\rho)\ln2}<\frac{25}{24},
\]
using its series, $\rho\le7/25$ and $\ln2>2/3$. Thus the bracket increases on both intervals. The accompanying script \texttt{support/network\_bounds.py} checks exactly these five rational inequalities; it does not test the probabilistic or asymptotic arguments.

\subsection{From truncation to all representations}\label{net:completion}
\begin{theorem}[Certificate-network contraction]\label{net:main}
There is an absolute $c>0$ such that for every sufficiently large dyadic $n$ and every nontrivial irreducible of dimension $D$ not annihilated by a pair layer,
\[
 \|T_n(\lambda)\|_{\op}\le D^{-c}.
\]
Consequently a fixed absolute number of sweeps has total-variation distance tending to zero.
\end{theorem}
\begin{proof}
Write $k=n-\lambda_1$ and $\mu=\bar\lambda$. For $k\le.14n$ and $k\le r\le n-k$, \eqref{eq:tuplemultiplicity} gives $m=\binom rk f^\mu$ and $D\le\binom nk f^\mu$.
If $k/n$ is bounded below, take $r=2k$. Proposition~\ref{net:densetruncation} applies uniformly over this density interval. Then $\log_2m\ge r-o(r)+\log_2 f^\mu$. The dense cutoff with slack $.01$ is at most $.97r$ and loses exponentially small mass, so Lemma~\ref{net:truncation} gives a fixed exponential saving when $\log f^\mu=O(r)$. If $\log f^\mu$ is larger than a sufficiently large fixed multiple of $r$, use instead the cutoff $L=\frac12\log_2m$. Markov's inequality and \eqref{net:sparsemgf} give $p\le2^{O(r)-bL}$, which is a fixed inverse power of $m$. These alternatives give a fixed inverse power of $D$.

If $k/n$ is small, fix $0<\gamma<\delta/(1+\delta)$ and take $r=\lfloor k(n/k)^\gamma\rfloor$. Then $r(r/n)^\delta=O(k)$, whereas
\[
 \log m\ge c_\gamma k\log(n/k)+\log f^\mu.
\]
Choose the fixed density cutoff sufficiently small. At $L=\frac12\log_2m$, Markov's inequality gives discarded mass $\exp(O(k))m^{-b/2}$, an inverse power of $m$. The retained bound is $m^{-1/4}$. Again $\log m$ is a fixed positive fraction of $\log D$.

For the remaining shapes use all $n$ cards, whose multiplicity is $D$. Lemma~\ref{lem:annihilation} removes first columns longer than $n/2$. Every other remaining shape has first row at most $.86n$. Let $a$ be the larger of its first row and column, transposing for the dimension estimate if necessary. If $.14n\le a\le.86n$, the first-row hook inflation over $a!$ is $\exp(O((n/a)\log n))=\exp(o(n))$. To verify this, for columns $j\le a/2$ sum $\lambda'_j-1$ and divide by $a/2$; for later columns use $\lambda'_j\le n/j\le2n/a$ and the harmonic sum of $1/(a-j+1)$. The remaining hook product is the hook product of the tail, whose tableau dimension is at least one. Hence
\[
 D\ge\binom na\exp(-o(n)),\qquad \log_2D\ge(.58-o(1))n,
\]
since $H_2(.14)>.58$. If $a<.14n$, the mean hook length is at most $a$: indeed the sum of hooks is $\frac12(\sum_i\lambda_i^2+\sum_j(\lambda'_j)^2)\le an$. Arithmetic--geometric mean and Stirling give a stronger bound. The full-density cutoff below $.54n$, with small slack, therefore beats the multiplicity exponentially. If $\log D$ is larger than a sufficiently large fixed multiple of $n$, switch to $L=\frac12\log_2D$ and \eqref{net:sparsemgf} as before. This proves the stated power bound for $K_n$ and, after taking square roots, for $T_n$.

For a sufficiently large fixed integer $r$, Plancherel and $\|T_n(\lambda)^r\|_{\HS}^2\le D\|T_n(\lambda)\|_{\op}^{2r}$ bound the squared relative $L^2$ distance by $\sum_{\lambda\ne(n),(1^n)}D^{2-2rc}$. Choose $r$ so the exponent is at most $-1$ and apply Lemma~\ref{lem:degrees}. Alternatively, in the range $k\le.14n$, the more precise hook estimate~\eqref{eq:near-row-hooks} bounds the same sum directly level by level. Lemma~\ref{lem:support} gives the matching lower order in physical time.
\end{proof}
